% VI/44: Supplementary reader bridge. Original problems and proofs remain intact.
\paragraph{A birational map is not necessarily a morphism.}
The classical quadratic Cremona transformation is the rational
map
\[
 \Phi:\mathbb P^2_k\dashrightarrow\mathbb P^2_k,
 \qquad [X:Y:Z]\longmapsto[YZ:ZX:XY].
\]
All three output forms have degree two, so rescaling the input
by $\lambda$ rescales the output by $\lambda^2$ wherever the
output triple is not identically zero.  But projective coordinates
$[0:0:0]$ do not define a point.  Accordingly, the expression is
not defined where $YZ=ZX=XY=0$, namely at the three coordinate
vertices
\[
 P_X=[1:0:0],\qquad P_Y=[0:1:0],\qquad P_Z=[0:0:1].
\]
These are the \emph{base points}.  A homogeneous formula does
not automatically give an everywhere-defined morphism.

\paragraph{Worked example 1: the involution on the dense torus.}
On the dense open set $XYZ\ne0$, the same formula can be
applied twice.  Writing $[U:V:W]=[YZ:ZX:XY]$, we find
\[
 \begin{aligned}
 \Phi^2([X:Y:Z])
   &=[VW:WU:UV]\\
   &=[X^2YZ:XY^2Z:XYZ^2]\\
   &=[X:Y:Z].
 \end{aligned}
\]
The last equality cancels the common nonzero factor $XYZ$.
Thus $\Phi$ restricts to an involutive isomorphism of the
dense algebraic torus; it is a \emph{birational involution} of
$\mathbb P^2$, not a biregular automorphism of the whole plane.

\paragraph{Worked example 2: three lines are contracted.}
On $X=0$ away from the two base points, $Y$ and $Z$ are
nonzero, so $\Phi([0:Y:Z])=[1:0:0]=P_X$.
Similarly, the coordinate lines $Y=0$ and $Z=0$ are mapped
to $P_Y$ and $P_Z$, respectively, wherever $\Phi$ is defined.
A whole one-dimensional locus can therefore collapse to a
single point under a birational \emph{rational} map.  This
behavior prevents $\Phi$ from being an automorphism and is
not contradicted by its being invertible on a dense open subset.

\paragraph{Worked example 3: a general line becomes a conic.}
Let $L\subset\mathbb P^2$ be the line $X+Y+Z=0$, which
contains none of the three coordinate vertices.  On the dense
torus, the inverse Cremona formula is
$[X:Y:Z]=[VW:WU:UV]$.  Substitution into the line equation
gives
\[
 VW+WU+UV=0.
\]
This homogeneous quadratic equation defines the closure of
the image $\Phi(L)$: a projective conic in the target with
coordinates $[U:V:W]$.  Its being degree two reflects the
fact that the coordinate expressions of $\Phi$ have degree
two, provided the original curve passes through no base points.
The same computation works for any line
$\alpha X+\beta Y+\gamma Z=0$ whose coefficients are all
nonzero, giving $\alpha VW+\beta WU+\gamma UV=0$.

\paragraph{Worked example 4: a line through a base point.}
Now choose $L':X+Y=0$, which passes through $P_Z$.
Substitution gives
\[
 W(V+U)=0.
\]
This is the \emph{total inverse-image equation} in the target:
it has an extra factor $W$, corresponding to an exceptional
coordinate line.  The closure of the birational transform of
$L'$ is instead the line $U+V=0$.  Thus the strict transform
has degree one, not two.  Confusing total transforms with the
strict transform produces spurious components in explicit
Cremona calculations.

\paragraph{How degrees change, and why blow-ups appear next.}
For an irreducible plane curve $C$ of degree $d$ that is not
contracted, let $m_X,m_Y,m_Z$ be its multiplicities at the
three base points.  The degree of its birational transform under
the quadratic Cremona map is
\[
 d'=2d-m_X-m_Y-m_Z.
\]
For the line $L$ of the previous example all multiplicities
are zero, giving $d'=2$.  For $L'$ only $m_Z=1$, giving
$d'=1$.  One explains the formula geometrically by resolving
the three base points: on the blow-up the pullback of a target
line has divisor class $2H-E_X-E_Y-E_Z$.  Intersection with
the proper transform of $C$ gives the stated degree formula.
The construction of the exceptional divisors and this
intersection computation belong to the following blow-up
section, VI/45; the examples here motivate it.

\paragraph{What the rational map does not promise.}
Cremona equivalence is much weaker than isomorphism of
projective varieties.  Indeterminacy may occur at base points,
curves may be contracted, and divisor degrees may change.
A correct argument therefore states the domain of definition,
computes the image on a dense open set, and only then takes the
scheme-theoretic closure or resolves the map as appropriate.
